\documentclass[10pt,a4paper]{amsart}
\usepackage[margin=19mm]{geometry}
\usepackage{amsmath,amssymb,mathtools}
\usepackage[scr=rsfs,cal=euler]{mathalfa}
\usepackage{xfrac}
\usepackage[unicode,hidelinks,bookmarks=false]{hyperref}
\newcommand{\ie}{i.e.\ }
\newcommand{\cf}{cf.\ }
\newcommand{\resp}{respectively\ }
\newcommand{\Subsection}{\subsection}
\providecommand{\nobreakdash}{\nobreak\hbox{-}}
\setlength{\emergencystretch}{2em}
\multlinegap0pt
\usepackage{bm}
\usepackage{bbm}
\usepackage{dsfont}
\usepackage{xspace}
\usepackage{cancel}
\usepackage{mathtools}
\usepackage{amsthm}
\makeatletter
\@ifundefined{X}{\newtheorem{X}{X}}{}
\@ifundefined{lemma}{\newtheorem{lemma}[X]{Lemma}}{}
\@ifundefined{proposition}{\newtheorem{proposition}[X]{Proposition}}{}
\makeatother
\makeatletter
\newcommand{\cA}{\ensuremath{\mathcal{A}}}
\newcommand{\cB}{\ensuremath{\mathcal{B}}}
\newcommand{\cD}{\ensuremath{\mathcal{D}}}
\newcommand{\cI}{\ensuremath{\mathcal{I}}}
\newcommand{\cJ}{\ensuremath{\mathcal{J}}}
\newcommand{\cS}{\ensuremath{\mathcal{S}}}
\newcommand*{\defeq}{\mathrel{\vcenter{\baselineskip0.5ex \lineskiplimit0pt
			\hbox{\scriptsize.}\hbox{\scriptsize.}}}%
	=}
\newcommand*{\eqdef}{=\mathrel{\vcenter{\baselineskip0.5ex \lineskiplimit0pt
			\hbox{\scriptsize.}\hbox{\scriptsize.}}}%
}
\expandafter\ifx\csname sproof\endcsname\relax
\newenvironment{sproof}{%
	\renewcommand{\proofname}{Sketch of Proof}\proof}{\endproof}
\fi
\makeatother
\title[Piatetski-Shapiro Primes in short intervals]{Piatetski-Shapiro Primes in short intervals}
\author{Lingyu Guo, Victor Zhenyu Guo}
\date{Source: arXiv:2606.01115v1 (CC BY 4.0)}
\begin{document}
\maketitle

\noindent\emph{Source and licence.} Piatetski-Shapiro Primes in short intervals. Version arXiv:2606.01115v1 (CC BY 4.0), distributed under the Creative Commons Attribution 4.0 International licence, as recorded in the arXiv licence metadata for this exact version and shown on the abstract page https://arxiv.org/abs/2606.01115v1. Full rights notice: licenses/arxiv-2606.01115-CC-BY-4.0.txt.\par\smallskip

\noindent\textit{Editorial scope.} Counted unit: the Lemma whose source label is \texttt{The Fundamental Lemma} in the article (source0.tex lines 1349--1372, source label \texttt{The Fundamental Lemma}), delivered together with the article's own complete original proof of it (source0.tex lines 1374--1464, one \texttt{proof} environment of 91 lines). No mathematics is written, repaired or re-derived: statement and proof are the article's own text line for line. Required context retained uncounted: lem:type1 (lines 823--836); condition of type1 (lines 829--831); lem:type2\_2 (lines 838--851); condition of type2\_2 (lines 844--846); lem:type2\_1 (lines 853--866); condition of type2\_1 (lines 859--861); harman sieve type I II information (lines 1190--1217); harman sieve type I (lines 1196--1198); a\_mn (lines 1292--1310); definition of a\_mn (lines 1294--1297). Every definition, display and label of those lines is retained unchanged. Omitted from this package and not redistributed: 1--822, 832--837, 847--852, 862--1189, 1199--1291, 1298--1348, 1373--1373, 1465--1765. Nothing inside a retained excerpt is omitted or altered: every retained body is delivered exactly as the article's lines are written, so no omission receipt of any kind accompanies this package. Citations: the retained excerpts carry no citation command at all, so no bibliography entry of the article is transcribed and no reference is redistributed. Reference adaptation: none was needed and none was made; every reference inside the delivered unit resolves inside it, so no unresolved reference or citation is produced. Printed slips kept literal: no printed word, symbol, hypothesis or number of the counted statement or of its proof was altered. Layout and class: the article's own class is \texttt{amsart} with packages including amsmath, amssymb, amsbsy, amsfonts, amsthm, latexsym, amsopn, amstext, amsxtra, euscript, amscd, stmaryrd, mathrsfs, cite; the wrapper is the lane's pinned \texttt{amsart} editorial wrapper at 10pt on A4, which restates the theorem-like environments the retained text needs and adds the article's own macro definitions that the retained text needs (\texttt{\textbackslash cA}, \texttt{\textbackslash cB}, \texttt{\textbackslash cD}, \texttt{\textbackslash cI}, \texttt{\textbackslash cJ}, \texttt{\textbackslash cS}, \texttt{\textbackslash defeq}, \texttt{\textbackslash eqdef}), with extra wrapper packages (bm, bbm, dsfont, xspace, cancel, mathtools). The delivered numbering is the unit's own. Assets: no retained span contains a figure environment, a graphics-inclusion command, a PDF-page inclusion, a table or a TikZ picture; no image file is copied into this package, the package declares no asset dependency and no image file is redistributed with it. Licence: version arXiv:2606.01115v1 of this article is distributed under the Creative Commons Attribution 4.0 International licence, as recorded in the arXiv licence metadata for this exact version and shown on the abstract page https://arxiv.org/abs/2606.01115v1. A full rights notice is delivered at licenses/arxiv-2606.01115-CC-BY-4.0.txt. Characters: every retained excerpt is pure ASCII, so the delivered engine, XeLaTeX with its default Latin Modern text font, reports no missing character.
\begin{proposition}[Type I]\label{lem:type1}
Let $\theta \in (2/3,1]$ and $N$ satisfies the condition
\begin{align*}
N \ll x^{2\gamma -\frac{4}{3}-\varepsilon} ,
\end{align*}
and
\begin{align}\label{condition of type1}
\frac{8-3\theta}{6} < \gamma <1,
\end{align}
then
$$
S_I \ll x^{\theta-\varepsilon}.
$$
\end{proposition}

\begin{proposition}[Type II]\label{lem:type2_2}
	Let $\theta \in (2/3,1]$ and $N$ satisfies the condition
	\begin{align*}
		x^{3-\theta-2\gamma+\varepsilon} \ll N \ll x^{\theta + 4\gamma -4-\varepsilon} \quad \text{or} \quad x^{5-\theta-4\gamma +\varepsilon}\ll N \ll x^{\theta+2\gamma-2-\varepsilon},
	\end{align*}
	and
	\begin{align}\label{condition of type2_2}
		\max\Big(\frac{1+\theta}{2} , \frac{6-3\theta}{4}\Big) < \gamma <1,
	\end{align}
	then
	$$
	S_{II} \ll x^{\theta-\varepsilon}.
	$$
\end{proposition}

\begin{proposition}[Type II]\label{lem:type2_1}
	Let $\theta \in (3/4,1]$ and $N$ satisfies the condition
	\begin{align*}
		x^{2-\theta-\gamma +\varepsilon}\ll N \ll x^{3\theta+5\gamma-7-\varepsilon} \quad \text{or} \quad x^{8-3\theta-5\gamma+\varepsilon} \ll N \ll x^{\theta + \gamma -1-\varepsilon} ,
	\end{align*}
	and
	\begin{align}\label{condition of type2_1}
		\frac{9-4\theta}{6} < \gamma <1,
	\end{align}
	then
	$$
	S_{II} \ll x^{\theta-\varepsilon}.
	$$
\end{proposition}

\begin{lemma}\label{harman sieve type I II information}
Let $\varepsilon>0$ be small, $2/3 < \theta < 1$ and $a_n,b_m \ll x^{\varepsilon}$. For
$$
N \ll x^{2\gamma -\frac{4}{3} -\varepsilon},
$$
we have 
\begin{align}\label{harman sieve type I}
\sum_{n \sim N} \sum_{\substack{m \\ mn \in \cA}} a_n  = \sum_{n \sim N} \sum_{\substack{m \\ mn \in \cB}} a_n \gamma (mn)^{\gamma-1} + O(x^{\theta + \gamma - 1 -\varepsilon}).
\end{align}
provided that \eqref{condition of type1} holds.
Moreover, if
$$
x^{5-\theta-4\gamma +\varepsilon}\ll N \ll x^{\theta+2\gamma-2-\varepsilon} \quad \text{or} \quad x^{3-\theta-2\gamma+\varepsilon} \ll N \ll x^{\theta + 4\gamma -4-\varepsilon},
$$
then
\begin{align}\label{harman sieve type II_2}
	\sum_{n \sim N} \sum_{\substack{m \\ mn \in \cA}} a_n b_m  = \sum_{n \sim N} \sum_{\substack{m \\ mn \in \cB}} a_n b_m \gamma (mn)^{\gamma-1} + O(x^{\theta + \gamma - 1 - \varepsilon}),
\end{align}
provided that \eqref{condition of type2_2} holds. For $3/4 < \theta <1$ and
$$
x^{2-\theta-\gamma +\varepsilon}\ll N \ll x^{3\theta+5\gamma-7-\varepsilon} \quad \text{or} \quad x^{8-3\theta-5\gamma+\varepsilon} \ll N \ll x^{\theta + \gamma -1-\varepsilon},
$$ 
we still have
\begin{align}\label{harman sieve type II_1}
	\sum_{n \sim N} \sum_{\substack{m \\ mn \in \cA}} a_n b_m  = \sum_{n \sim N} \sum_{\substack{m \\ mn \in \cB}} a_n b_m \gamma (mn)^{\gamma-1} + O(x^{\theta + \gamma - 1 - \varepsilon}),
\end{align}
provided that \eqref{condition of type2_1} holds.
\end{lemma}

\begin{lemma}\label{a_mn}
Let $I$, $J$ be integers and $\cI_i$, $\cJ_j$ be intervals for $1 \leqslant i \leqslant I$, $1 \leqslant j \leqslant J$. Write
\begin{align}\label{definition of a_mn}
a_{mn} \defeq \sum_{\substack{kp_1 \cdots p_l = n \\ p_1 < p_2 < \cdots < p_l \\ p_i \in \cI_i}} c_{n}
\sum_{\substack{lq_1 \cdots q_j = m \\ q_1 < q_2 < \cdots < q_j \\ q_j \in \cJ_j}} d_{m}
\end{align}
with $|c_n|$, $|d_m| \leqslant x^{\varepsilon}$ and $p_1, \ldots, p_l$ and $q_1, \ldots, q_j$ satisfying $t$ joint conditions of the form
$$
p_u \leqslant q_v \quad \text{or} \quad q_v \leqslant p_u
$$
or
$$
\prod_{u \in U} p_u\prod_{v \in V} q_v \leqslant H \quad \text{or} \quad \prod_{u \in U} p_u \geqslant \prod_{v \in V} q_v
$$
or similar (for given $U \subset \{1, \ldots, I \}$, $V \subset \{1, \ldots, J\}$ and $H\leqslant x$). Suppose that $\varepsilon, \gamma, \theta, N$ satisfy the conditions of Proposition~\ref{lem:type2_2} or Proposition~\ref{lem:type2_1}, then
$$
\sum_{n \sim N} \sum_{\substack{m \\ mn \in \cA}} a_{mn} = \sum_{n \sim N} \sum_{\substack{m \\ mn \in \cB}} a_{mn} \gamma(mn)^{\gamma-1} + O(x^{\theta + \gamma - 1 -\varepsilon}).
$$
\end{lemma}

\begin{lemma}[Fundamental Lemma]\label{The Fundamental Lemma}
Let $u \geqslant 1$ be an integer, $\varepsilon > 0$ be a small positive number. 
	
(i) Let $p_1, \ldots, p_u \in [1, x]$ be such that $\prod_{1 \leqslant k \leqslant u} p_k \leqslant x$. Suppose that $\varepsilon, \gamma, \theta, N$ satisfy the conditions of Proposition~\ref{lem:type2_2} or Proposition~\ref{lem:type2_1} and there exists $\cD \subset \{1, \ldots, u\}$ such that  
$$\prod_{k \in D} p_k \asymp N.$$
Then  
\begin{equation}\label{fundamental lemma 1}
\sum_{\substack{ p_1 , \cdots , p_u  }} \cS(\cA_{p_1, \cdots, p_u}, p_1) = \gamma x^{\gamma-1} \sum_{\substack{p_1 , \cdots , p_u }} \cS(\cB_{p_1, \cdots, p_u}, p_1) + O(x^{\theta + \gamma - 1 - \varepsilon}).
\end{equation}

(ii)For any $a_n \ll x^{\varepsilon}$, if $\varepsilon, \gamma, \theta$ satisfy the conditions of Proposition~\ref{lem:type1} and satisfy the conditions of Proposition~\ref{lem:type2_1} or Proposition~\ref{lem:type2_2} we have
\begin{equation}\label{fundamental lemma 2}
\sum_{n \sim N} a_n \cS(\cA_n, z) = \gamma x^{\gamma-1} \sum_{n \sim N} a_n \cS(\cB_n, z) + O(x^{\theta + \gamma - 1 - \varepsilon} ),
\end{equation}
where 
$$
N \ll x^{3\theta + 5\gamma -7 -\varepsilon} \quad \text{and} \quad z = x^{4\theta+6\gamma-9-\varepsilon}
$$ 
for $\varepsilon, \gamma, \theta$ satisfying the conditions of Proposition~\ref{lem:type1} and \ref{lem:type2_1}, while 
$$
N \ll x^{\theta + 4\gamma -4-\varepsilon} \quad \text{and} \quad z = x^{2\theta+6\gamma-7-\varepsilon}
$$ 
when $\varepsilon, \gamma, \theta$ satisfy the conditions of Proposition~\ref{lem:type1} and Proposition~\ref{lem:type2_2}.
\end{lemma}

\begin{proof}
We prove \eqref{fundamental lemma 1} first. By the definition of $\cS(\cA_n,z)$, we have
\begin{align*}
\sum_{\substack{p_1 , \cdots , p_u}} \cS(\cA_{p_1, \cdots, p_u}, p_1) = \sum_{\substack{p_1 , \cdots , p_u}} \sum_{\substack{s \\ p_1\cdots p_u s \in \cA \\ (s, P(p_1))=1}} 1.
\end{align*}
Let 
$$
n = \prod_{j\in \cD} p_j, \quad m = s\prod_{j \notin \cD} p_j.
$$
Hence we arrive at
\begin{align} \label{fundamental lemma 3}
\nonumber
&\sum_{\substack{p_1 , \cdots , p_u}} \cS(\cA_{p_1, \cdots, p_u}, p_1) \\
\nonumber
& = \sum_{n \sim N} \sum_{\substack{m \\ mn \in \cA \\ (m,P(p_1)) = 1}} \bigg( \sum_{\substack{\prod_{j\in \cD} p_j = n}} 1 \bigg) \bigg( \sum_{ s\prod_{j \notin \cD} p_j = m} 1 \bigg) \\
\nonumber
& = \sum_{n \sim N} \sum_{\substack{m \\ mn \in \cA }} \bigg( \sum_{\substack{\prod_{j\in \cD} p_j = n}} 1 \bigg) \bigg( \sum_{ s\prod_{j \notin \cD} p_j = m} \sum_{d|(m,P(p_1))} \mu(d) \bigg)\\
& \eqdef \sum_{n \sim N} \sum_{\substack{m \\ mn \in \cA }} a_{mn},
\end{align}
where we use the identity
$$
\textbf{1}_{(m,n)=1} = \sum_{d|(m,n)} \mu(d) .
$$
Noting that $a_{mn}$ in \eqref{fundamental lemma 3} satisfies the definition of \eqref{definition of a_mn}, we apply Lemma~\ref{a_mn} to \eqref{fundamental lemma 3} and find that
\begin{align}\label{fundamental lemma 4}
\sum_{\substack{p_1 , \cdots , p_u}} \cS(\cA_{p_1, \cdots, p_u}, p_1) = \sum_{n \sim N} \sum_{\substack{m \\ mn \in \cB}} a_{mn} \gamma(mn)^{\gamma-1} + O(x^{\theta + \gamma - 1 -\varepsilon}).
\end{align} 
The proof is finished if we prove that the double sum of the right-hand side in \eqref{fundamental lemma 4} is
$$
\gamma x^{\gamma-1} \sum_{\substack{p_1 , \cdots , p_u}} \cS(\cB_{p_1, \cdots, p_u}, p_1) + O(x^{\theta + \gamma - 1 - \varepsilon}).
$$ 
Since $a_{mn} \geqslant 0$, we subtract the two main terms and use the trivial bound to get
\begin{align*}
\nonumber
&\sum_{n \sim N} \sum_{\substack{m \\ mn \in \cB}} a_{mn} \gamma(mn)^{\gamma-1} - \gamma x^{\gamma-1} \sum_{\substack{p_1 , \cdots , p_u}} \cS(\cB_{p_1, \cdots, p_u}, p_1) \\ 
\nonumber
=&\gamma \sum_{n \sim N} \sum_{\substack{m \\ mn \in \cB}} a_{mn} ((mn)^{\gamma-1}-x^{\gamma-1}) \\
\leqslant&  \gamma ( x^{\gamma-1} - (x+y)^{\gamma-1})\sum_{n \sim N} \sum_{\substack{m \\ mn \in \cB}} a_{mn} \ll y^2x^{\gamma-2} \ll x^{\theta +\gamma -1 - \varepsilon}.
\end{align*} 

As for \eqref{fundamental lemma 2}, we take $z = x^{4\theta + 6\gamma - 9 - \varepsilon}$ and suppose that $\varepsilon, \gamma, \theta, N$ satisfy the conditions of Proposition~\ref{lem:type1} and \ref{lem:type2_1}. We note that
$$
\sum_{n \sim N} a_n \cS(\cA_n, z) = \sum_{n \sim N} a_n \sum_{\substack{ m \\ mn\in \cA \\ (m,P(z))=1}}1 =\sum_{n \sim N} a_n \sum_{\substack{ d|P(z) \\ mnd \in \cA}} \mu(d),
$$
where we change the variable $m$ by $md$ in the last equality. Once we show that
\begin{align}\label{fundamental lemma 5}
\sum_{n \sim N} a_n \sum_{\substack{ d|P(z) \\ mnd \in \cA}} \mu(d) = \sum_{n \sim N} a_n \sum_{\substack{ d|P(z) \\ mnd \in \cB}} \mu(d) \gamma (mnd)^{\gamma-1} + O(x^{\theta + \gamma - 1 - \varepsilon}),
\end{align}
by the similar arguments of \eqref{fundamental lemma 1}, the proof is done. We divide the sum in the left-hand side of \eqref{fundamental lemma 5} into two parts
\begin{align*}
\Sigma_1 \defeq \sum_{n \sim N} a_n \sum_{\substack{ d|P(z) \\ mnd \in \cA \\ nd \leqslant x^{3\theta+5\gamma-7-\varepsilon}}} \mu(d), \quad
\Sigma_2 \defeq \sum_{n \sim N} a_n \sum_{\substack{ d|P(z) \\ mnd \in \cA \\ nd > x^{3\theta+5\gamma-7-\varepsilon}}} \mu(d).
\end{align*}
In $\Sigma_1$ we produce a new variable $k = nd$ and get
$$
\Sigma_1 = \sum_{k \leqslant x^{3\theta+5\gamma-7-\varepsilon}} \sum_{\substack{m \\ km \in \cA }} b_k, \quad \text{where} \quad b_k \defeq \sum_{\substack{nd=k \\ n \sim N \\ d|P(z)}}a_n \mu(d).
$$
Obviously $|b_k| \ll x^{\varepsilon}$, so $\Sigma_1$ is a type I sum. By \eqref{harman sieve type I} in Lemma~\ref{harman sieve type I II information}, we have that
$$
\Sigma_1 = \sum_{n \sim N} a_n \sum_{\substack{ d|P(z) \\ mnd \in \cB \\ nd \leqslant x^{3\theta+5\gamma-7-\varepsilon}}} \mu(d) \gamma (mnd)^{\gamma-1} + O(x^{\theta + \gamma - 1 - \varepsilon}).
$$
Now we write $p_1$ as the largest prime factor of $d$ and replace $d$ by $pd$ to obtain that
\begin{align}\label{fundamental lemma 6}
\Sigma_2 = \sum_{n \sim N} a_n \sum_{p_1 < z} \sum_{\substack{p_1d|P(z) \\ mnp_1d \in \cA \\ np_1d > x^{3\theta+5\gamma-7-\varepsilon}}} \mu(p_1d) = - \sum_{n \sim N}a_n \sum_{p_1 < z} \sum_{\substack{d|P(p_1) \\ mnp_1d \in \cA \\ np_1d > x^{3\theta+5\gamma-7-\varepsilon}}} \mu(d).
\end{align}
We divide the last sum of \eqref{fundamental lemma 6} into two parts , say
\begin{align*}
&\Sigma_3 \defeq \sum_{n \sim N}a_n \sum_{p_1 < z} \sum_{\substack{d|P(p_1) \\ mnp_1d \in \cA \\ nd \leqslant x^{3\theta+5\gamma-7-\varepsilon} < np_1d}} \mu(d) , \\
&\Sigma_4 \defeq \sum_{n \sim N}a_n \sum_{p_1 < z} \sum_{\substack{d|P(p_1) \\ mnp_1d \in \cA \\ nd > x^{3\theta+5\gamma-7-\varepsilon}}} \mu(d).
\end{align*}
As for $\Sigma_3$, we draw two new variables, $k = nd$ and $l=mp_1$. Since $p_1 < z$, we note that $p_1k > x^{3\theta+5\gamma-7-\varepsilon} \Rightarrow k > x^{2 - \theta - \gamma +\varepsilon}.$ Hence similar to the arguments of \eqref{fundamental lemma 3}, we have that $\Sigma_3$ is a type II sum in the form of 
$$
\Sigma_3 = \sum_{\substack{ x^{2-\theta -\gamma + \varepsilon} \leqslant k \leqslant x^{3\theta+5\gamma-7-\varepsilon}}} \sum_{\substack{ l \\ kl\in \cA}} a_{kl},
$$ 
where 
$$
a_{kl} \defeq  \sum_{\substack{ nd = k \\ n \sim N \\ d|P(p_1)}} a_n \mu(d)\sum_{\substack{mp_1 = l \\ p_1 < z}}1.
$$
By Lemma~\ref{a_mn}, we have
$$
\Sigma_3 = \sum_{n \sim N}a_n \sum_{p_1 < z} \sum_{\substack{d|P(p_1) \\ mnp_1d \in \cB \\ nd \leqslant x^{3\theta+5\gamma-7-\varepsilon} < np_1d}} \mu(d) \gamma (mnp_1d)^{\gamma-1} + O(x^{\theta + \gamma - 1 - \varepsilon}).
$$
We treat $\Sigma_4$ similarly to $\Sigma_2$ and write it as the sum of the following two sums
\begin{align*}
&\Sigma_5 \defeq \sum_{n \sim N}a_n \sum_{p_2 < p_1 < z} \sum_{\substack{d|P(p_2) \\ mnp_2p_1d \in \cA \\ nd \leqslant x^{3\theta+5\gamma-7-\varepsilon} < np_2d}} \mu(d) , \\
&\Sigma_6 \defeq \sum_{n \sim N}a_n \sum_{p_2 < p_1 < z} \sum_{\substack{d|P(p_2) \\ mnp_2p_1d \in \cA \\ nd > x^{3\theta+5\gamma-7-\varepsilon}}} \mu(d).
\end{align*}
We deal with $\Sigma_5$ as we did with $\Sigma_3$ to obtain a similar asymptotic formula and we give further decomposition for $\Sigma_6$. We can continue in this treatment to obtain each $\Sigma_{2j-1}$ for which we can apply Lemma~\ref{a_mn} and $\Sigma_{2j}$ for which we give further decomposition. Since the integers in the interval $(x, x+y]$ have $< \log x$ prime divisors after at most $\log x$ such steps we will obtain an empty $\Sigma_{2j}$. Thus we have given asymptotic formula for all the occurring sums. Clearly, combining the asymptotic formula for all $\Sigma_{2j-1}$ we complete the proof of \eqref{fundamental lemma 5} for $N \leqslant x^{3\theta +5\gamma-7 -\varepsilon}$.

When  $\varepsilon, \gamma, \theta, N$ satisfy the conditions of Proposition~\ref{lem:type1} and \ref{lem:type2_2}, the proof is the same. 
\end{proof}

\end{document}
